\usepackage{comment}
\usepackage{subfigure}
\usepackage{amsmath}
\usepackage{bm} % added by Harashima (2020/01/15)
\usepackage{ascmac} % added by Harashima (2020/02/16)
\usepackage{multirow} % added by Harashima (2020/03/09

\usepackage{booktabs}
\usepackage{bbding}
\usepackage{pifont}
\usepackage{wasysym}
\usepackage{amssymb}

\newcommand{\mytitle}{WeaveNet for Approximating Two-sided Matching Problems}


\makeatletter
\def\WordCount#1{%
\@tempcnta\z@%
\@tfor \@tempa:=#1\do{\advance\@tempcnta\@ne}%
\@tempcnta%
}%
\makeatother

\newif\ifprintjp

\newcommand{\para}[2]{%

\ifprintjp%
  \ifx"#1"% #1が空なら英語を表示しておく．
  #2
  \else%
  #1 
  \fi
\else%
  \ifx"#2"% #2が空なら日本語を表示しておく．
  #1
  \else%
  #2 
  \fi
\fi

}


\usepackage[dvipsnames]{xcolor}

%\newcommand{\writtenby}[1]{\textcolor{red}{Written by #1 and waiting for a proofread.}}
\newcommand{\writtenby}[1]{}
%\newcommand{\checkedby}[2]{\textcolor{blue}{Proofread by #1 (originally written by #2).}} 
\newcommand{\checkedby}[2]{}
%If you critically modify something in this paragraph, please remove this \textbackslash checkedby flag.}}

\newcommand{\modifiedby}[2]{\textcolor{magenta}{Modified by #1 (Originally written by #2)}}

\usepackage{url}
% 各人のコメント用コマンド．
%\begin{comment}
\newcommand{\ah}[1]{\textcolor{BurntOrange}{[AH: #1]}}
\newcommand{\sone}[1]{\textcolor{RubineRed}{[SS: #1]}} %ssは既に存在しているので，曾根さんだけ\ssではなく\soneでお願いします．
\newcommand{\jm}[1]{\textcolor{ForestGreen}{[JM: #1]}}
%\end{comment}

% コメントをそのまま採用する場合は，上をコメントアウトし，ここをコメントインする．
\begin{comment}
\newcommand{\ah}[1]{#1}
\newcommand{\sone}[1]{#1}
\newcommand{\jm}[1]{#1}
\end{comment}


% math
\newcommand{\loss}{{\mathcal L}}
\newcommand{\lowest}{C_{min}}
\newcommand{\ull}[1][\underline]{#1}
\newcommand{\oll}[1][\overline]{#1}

\newcommand{\SEq}{S\hspace{-0.1em}Eq}
